\section{Undirected connectivity: exact sums, optimization, and witnesses}
\label{sec:undirected-proof}

We prove Theorem~\ref{thm:undirected}, always assuming \HTC.
Let $S_G=n+m+2$. The determinant scale in
Section~\ref{sec:determinants} is
\[
 T_{\mathrm{det}}=n+(n+2m)+2=2n+2m+2=O(S_G).
\]
These are distinct parameters; the precision below is measured in
$S_G$.

\subsection{Local candidate systems}

Write $u_v$ for selected-vertex bits and $x_e$ for selected-edge
bits. Requiring selected-edge endpoints to be selected is a constant
scope condition. Incidence degree conditions, including exact degrees
zero, one, and two, have $O(\deg_G(v)+1)$-wire functional chains:
pass the saturated count $\min\{3,\text{count so far}\}$ in two
bits. Domination is checked by an OR chain on each closed neighborhood,
with a total of $O(n+m)$ gates. Vertex cover has the binary constraints
$u_v\vee u_w$ for every edge $vw$. In induced systems introduce the
unique bit $x_{vw}=u_vu_w$ by an AND relation. All such systems have
$O(S_G)$ original and auxiliary bits and total scopes $O(S_G)$.
Global cardinalities and cost relations are the fixed number of
statistics from Section~\ref{subsec:weights}.

For each candidate choose the matrix~\eqref{eq:marker-matrix} and
combine its sparse determinant gadget with these constraints. The
contraction is exactly
\begin{equation}\label{eq:object-sum}
 Z=\sum_{y\text{ satisfying the local and target conditions}}\det A(y).
\end{equation}
The sum over permutation bits is taken inside the sum over $y$; all
other auxiliary states have unique extensions. Table~\ref{tab:objects}
lists the candidates and the determinant value on the objects that
survive the filter. Each single-root row explicitly requires $r\in U$.

\begin{table}[ht]
\centering\small
\renewcommand{\arraystretch}{1.16}
\begin{tabular}{@{}P{.23\linewidth}P{.44\linewidth}P{.25\linewidth}@{}}
\toprule
Problem & Candidate constraints and markers & Surviving weight\\
\midrule
HC / Hamiltonian TSP & $U=V$, all selected degrees two; $M_0=\{r\}$ & $n$ for a Hamiltonian edge set\\
Hamiltonian path & $U=V$, degrees one at distinct $s,t$, two elsewhere; $M_0=\{s\}$ & $1$ for a spanning $s$--$t$ path\\
Simple path & $s,t\in U$, degree one at $s,t$, two at other selected vertices, zero outside; $M_0=\{s\}$ & $1$ for an $s$--$t$ path on $U$\\
Steiner tree & Terminals in $U$, $F\subseteq E_G[U]$, $|F|=|U|-1$; $M_0=\{r\}$ for a fixed terminal & $1$ for a tree $(U,F)$\\
Connected vertex cover & Cover constraints, $F=E_G[U]$, $M_0=\{r\}$ & $\tau(G[U])>0$\\
Connected domination & Domination constraints, $F=E_G[U]$, $M_0=\{r\}$ & $\tau(G[U])>0$\\
Feedback vertex set & $F=E_G[U]$, $M_0\subseteq U$, $|F|+|M_0|=|U|$ & $1$ for each forest with one marker per component\\
\bottomrule
\end{tabular}
\caption{In every row the determinant is zero on candidates failing
the required connectivity condition. Costs and a specified path length
can be added as target relations.}
\label{tab:objects}
\end{table}

Here is the verification of each row. A simple graph of degree two at
every vertex is a disjoint union of cycles. A single marker allows
only the component containing that marker. A surviving cycle on $n$
vertices has exactly $n$ spanning trees, obtained by deleting one
edge. Consequently
\begin{equation}\label{eq:hc-readback}
 Z=n\,\#\mathrm{HC}.
\end{equation}
No orientation or starting-vertex multiplicity is present: the objects
are undirected edge sets. The matrix remains unweighted even when
cost constraints are imposed on these edge sets.

For path candidates, the two odd-degree endpoints must lie in the
same path component, and every other component is a cycle. The
single marker eliminates those cycles; the remaining path has tree
number one. For a specified length $\ell$, impose $|U|=\ell+1$.
For Steiner candidates, connectivity and $|F|=|U|-1$ make the chosen
subgraph a tree. The counted objects are all tree subgraphs containing
the terminal set, including ones with nonterminal leaves. For
connected covers and dominating sets, the determinant is their induced
subgraph's positive tree count, so positivity of~\eqref{eq:object-sum}
is exactly existence.

For FVS, append a root adjacent to every marker. When the determinant
is positive, this augmented graph is connected, and the imposed equation
makes its number of edges equal to its vertex count minus one. It is
therefore a tree and the determinant equals one. Removing its root
leaves the induced forest $G[U]$, with exactly one marker per component.
Conversely every induced forest has such marker choices, all of weight
one. Thus $Z>0$ decides whether a forest of the specified deletion
cost exists. Different marker choices give a positive multiplicity
for the same deletion set.

\subsection{Precision and exact readback}

All matrices in~\eqref{eq:object-sum} have nonnegative determinant by
\eqref{eq:matrix-tree}. Each augmented graph has at most $m+n$ edges,
so its tree count is at most $2^{m+n}$. The numbers of candidates
$(U,F)$, $(U,M_0)$, or $U$ are at most $2^{n+m}$, $3^n$, and $2^n$,
respectively. Uniformly,
\begin{equation}\label{eq:graph-readback-bound}
 0\le Z\le2^{\max\{2n+2m,\,3n+m\}}<2^{4S_G}.
\end{equation}
Take
\begin{equation}\label{eq:graph-precision}
 P_0=4S_G+1,\qquad \ell=1.
\end{equation}
The local signed entries, including $-1$, are legitimate residues in
$\Z/2^{P_0}\Z$. Their intermediate products may cancel; the final
sum is the bounded nonnegative integer in
\eqref{eq:graph-readback-bound}. Its least nonnegative residue is
therefore exactly $Z$. For HC divide that integer by $n$, as in
\eqref{eq:hc-readback}. Even when $n$ is even, no inverse modulo a
power of two is required.

The local candidate constraints, degree counters, and numerical
matrix factors have linear total scope length. The sparse sign gadget
has degree $O(S_G^{2/3})$; adding long cost constraints raises the
degree to $O(S_G^{9/10})$ and keeps $O(S_G)$ wires. Choose
$s=\lceil K S_G\rceil$ with a fixed sufficiently large $K$, and
subdivide edges for bipartiteness. Precision~\eqref{eq:graph-precision},
vertex count, description length, and polynomial-time local evaluation
then satisfy every part of \HTC. One query takes
\begin{equation}\label{eq:one-graph-query}
 \poly(S_G+L)2^{O(g(S_G))}
\end{equation}
time and polynomial space.

\subsection{Optimization, witnesses, and boundary conventions}

Nonnegative edge and vertex costs are attached to the \emph{object}
bits, not to the spanning trees counted by the determinant. For FVS
the objective is the cost of $V\setminus U$, implemented with
complement bits. Use a cost-at-most predicate for minimization and
a cost-at-least predicate for maximization. The sum of the input
costs has $O(n^{9/10}+\log n)$ bits on a sparse graph, and binary
search needs only polynomially many queries. For longest path one
may maximize the number of selected vertices, or an explicitly given
nonnegative selected-edge/vertex weight. The long-coefficient gadget
handles these exact target relations without enumerating cost values.

If a target has $Z>0$, fix its object bits successively. For a bit $y_i$
we have $Z=Z_{y_i=0}+Z_{y_i=1}$, with both summands nonnegative.
Keep any positive child. Apply this only to object bits (and, for
FVS, marker bits), retaining the determinant's permutation sum.
At the end the surviving object satisfies the target and its graph
property. Direct polynomial-time graph and cost checks verify the
returned witness. At most a linear number of additional queries is
needed.

Enumerating possible roots and endpoint pairs costs only a polynomial
factor. For exact counting of unspecified-endpoint undirected paths,
use $s<t$ to count each nontrivial edge set once. Single-vertex paths
are counted separately. For prescribed $s=t$, only the zero-length
single-vertex path is included. Simple undirected cycles require at
least three vertices. An empty terminal set allows the empty Steiner
tree of cost zero in optimization. For counting all nonempty Steiner
tree subgraphs with no terminals, choose the unique root $r=\min U$
and force smaller vertices outside $U$; add the empty tree only if
that convention is requested. With a nonempty terminal set, fix one
terminal as the root, so that no root multiplicity is introduced.
Single-root trees are included. Empty-cover and empty-graph conventions
are checked directly; all nonempty connected sets are covered by the
root enumeration.

The number of queries, including roots, endpoints, binary search, and
self-reduction, is polynomial. Their workspace is reused. Since
$m\le cn$, $g(S_G)=O_c(g(n))$, proving
Theorem~\ref{thm:undirected}.
